\chapter{Reproducing the Prototype and its Evaluation}
\label{app:reproduce}

This appendix lists the steps needed to run the prototype and repeat the measurements of \Cref{ch:results}. The source code is in the repository named in the project README; the commands assume Linux or WSL2 with Docker, Node.js 20.9 or newer, and the Hyperledger Fabric binaries and \texttt{test-network} of the \texttt{fabric-samples} repository.

\section{Starting the Network and the Application}
\label{app:start}

\begin{lstlisting}[numbers=none, language={}, basicstyle=\ttfamily\footnotesize]
cd ~/fabric-samples/test-network
./network.sh down
./network.sh up createChannel -c mychannel
./network.sh deployCC -ccn pmcc -ccp ../dpm-project/pm-chaincode \
    -ccl javascript -ccv 1.0 -ccs 1

cd ~/fabric-samples/dpm-project/pm-backend  && npm install && npm run dev
cd ~/fabric-samples/dpm-project/pm-frontend && npm install && npm run dev
\end{lstlisting}

The chaincode functions changed signature when permissions were added, so after any chaincode change it is simplest to recreate the network with \texttt{./network.sh down}. Demonstration data (five accounts, four projects and their tasks) can be loaded with \texttt{npm run seed} in the backend folder.

\section{Tests and Benchmark}
\label{app:tests}

\begin{lstlisting}[numbers=none, language={}, basicstyle=\ttfamily\footnotesize]
cd pm-backend  && npm test            # backend suite, no network needed
cd pm-frontend && npm test            # frontend decision logic
cd pm-backend  && npm run benchmark   # needs the running network and backend
npm run verify-audit -- audit-task-<id>.json   # offline check of a report
\end{lstlisting}

The benchmark writes its measurements to \texttt{benchmark-results.json} and \texttt{benchmark-results.csv}, from which \Cref{tab:benchmark} is filled in.

\section{Summary of the REST API}
\label{app:api}

\Cref{tab:api} lists the main endpoints. All endpoints other than sign-in and sign-up require a valid session cookie, and state-changing requests are checked against the \texttt{Origin} and \texttt{Host} headers.

\begin{table}[tb]
    \centering
    \small
    \begin{tabular}{p{1.3cm}p{5.6cm}p{5.4cm}}
        \toprule
        Method & Path & Chaincode function \\
        \midrule
        POST & \texttt{/api/projects} & \texttt{createProject} \\
        POST & \texttt{/api/projects/\{id\}/members} & \texttt{addProjectMember} \\
        DELETE & \texttt{/api/projects/\{id\}} & \texttt{archiveProject} \\
        GET & \texttt{/api/projects/\{id\}/tasks} & \texttt{getProjectTasks} \\
        GET & \texttt{/api/projects/\{id\}/audit} & \texttt{getProjectHistory} + ledger check \\
        POST & \texttt{/api/tasks} & \texttt{createTask} \\
        PUT & \texttt{/api/tasks/\{id\}/assign} & \texttt{assignTask} \\
        PUT & \texttt{/api/tasks/\{id\}/status} & \texttt{updateTaskStatus} \\
        PUT & \texttt{/api/tasks/\{id\}/meta} & \texttt{updateTaskMeta} \\
        POST & \texttt{/api/tasks/\{id\}/comments} & \texttt{addComment} \\
        POST & \texttt{/api/tasks/\{id\}/files} & \texttt{attachFile} \\
        DELETE & \texttt{/api/tasks/\{id\}} & \texttt{archiveTask} \\
        GET & \texttt{/api/tasks/\{id\}/audit} & \texttt{getTaskHistory} + ledger check \\
        \bottomrule
    \end{tabular}
    \caption{Main REST endpoints and the chaincode functions they invoke.}
    \label{tab:api}
\end{table}
